% Source macros retain the measured results and provenance in numbers_results.tex.
% Split the former 12-column table for readable reconstruction/cost comparisons.
% 12 Nov revision: descriptive model names (Final = teacher, Control = uncorrected teacher, Base = initial model);
% protocol names (V1/V0, v23) removed from the heads; caveats moved into the note below the table and into
% Section 5.2 (sec:res-fairness).
\begin{table}[!tbp]
\centering\small
\setlength{\tabcolsep}{3pt}
\begin{tabular}{@{}lrrrrrrc@{}}
\toprule
& \multicolumn{3}{c}{AnnyOne holdout} & \multicolumn{3}{c}{3DPW test} & \\
\cmidrule(lr){2-4}\cmidrule(lr){5-7}
Model & PVE & PA-PVE & Recall & PVE & PA-PVE & Recall & \shortstack{3DPW\\training} \\
\midrule
Released baseline & \NAVERAnnyOne$^a$ & \NAVERPAAnnyOne$^a$ & \NAVERRecallAnnyOne$^a$ & \NAVERThreeDPW & \NAVERPAThreeDPW & \NAVERRecallThreeDPW & no \\
Teacher & \FinalPVE & \FinalPAPVE & \FinalRecall & \textbf{\FinalThreeDPW} & \textbf{\FinalPAThreeDPW} & \FinalRecallThreeDPW & yes \\
Student & \StudentPVE & \StudentPAPVE & \StudentRecall & \StudentThreeDPW & \StudentPAThreeDPW & \StudentRecallThreeDPW & yes \\
Uncorrected teacher & \ControlPVE & \ControlPAPVE & \ControlRecall & \ControlThreeDPW & \ControlPAThreeDPW & 99 & yes \\
Initial model & \BasePVE & \BasePAPVE & \BaseRecall & \BaseThreeDPW & \BasePAThreeDPW & 100 & no \\
\bottomrule
\multicolumn{8}{@{}p{0.97\linewidth}@{}}{\scriptsize $^a$Potentially in-sample: the released baseline was trained
on the AnnyOne release, which contains the holdout images. 3DPW errors are measured against the converted Anny
labels and are not comparable with published SMPL-space results (Section~\ref{sec:data-limits}).}
\end{tabular}
\caption[Reconstruction of the student, the teacher and the released baseline]{Reconstruction of the student, its
teacher and the released baseline on the AnnyOne holdout and the 3DPW test set (errors in mm, recall in \%; one run
per model). The models differ in their training data, and the released baseline may have seen the holdout, so a
lower PVE does not show a better architecture (Section~\ref{sec:res-fairness}).}
\label{tab:baseline}
\end{table}

\begin{table}[!tbp]
\centering\small
\begin{tabular}{@{}lrrrr@{}}
\toprule
Model & Params (M) & GFLOPs & \shortstack{Latency\\L40S (ms)} & \shortstack{Latency\\A100 (ms)} \\
\midrule
Released baseline & \NAVERParams & $\geq$\NAVERBackboneGFLOPs$^b$ & n/m & n/m \\
Teacher & \FinalParams & \FinalGFLOPs & \TeacherMsLFortyS & \TeacherMsAHundred \\
Student & \StudentParams & \StudentGFLOPs & \StudentMsLFortyS & n/m \\
\bottomrule
\multicolumn{5}{@{}p{0.8\linewidth}@{}}{\scriptsize One multiply-add counts as two FLOPs. Latency: median batch-one
forward time in fp32 with TF32; the student and teacher L40S times come from the same job. n/m: not measured.
$^b$Backbone only; the total is higher. The initial model and the uncorrected teacher share the teacher
architecture, and the fixed rotation corrections add no trainable parameters.}
\end{tabular}
\caption[Model size, computation and forward latency]{Model cost at 672\,px input with one person. Latencies are
comparable only within a GPU column.}
\label{tab:model-cost}
\end{table}
